\section{Multiplicity transport and stable labels}
\label{sec:cluster-transport}

The current FIBER audit identifies two specific unresolved issues: a short
edge in an incidence graph does not give the same diameter bound for its
whole component, and a metric gap between complex roots does not stabilize
lexicographic ordering of rational boxes~\cite{repo-fiber}.  We resolve the
finite mathematical issues here.  A polynomial application must still
produce a sound complete multiplicity certificate.  No matching or
certificate-independent observation is assumed in the theorems below.

\subsection{The certificate interface}

Let $\Omega=\{0,\ldots,N-1\}$ be a finite set of \emph{occurrences}; repeated
roots have separate occurrences.  Let $x_t$, $t<N$, be points of a metric
space.  A weighted enclosure consists of centers $a_0,\ldots,a_{m-1}$,
positive natural weights $r_0,\ldots,r_{m-1}$, a radius $\epsilon_A\geq0$,
and a coverage witness
\[
 \alpha:\Omega\longrightarrow\{0,\ldots,m-1\},\qquad
 |\alpha^{-1}(i)|=r_i,\qquad d(x_t,a_{\alpha(t)})\leq\epsilon_A.
 \tag{C-A}
\]
A second enclosure $(b_j,s_j,\epsilon_B,\beta)$ has the analogous properties
for the \emph{same} occurrence multiset.  The occurrence data may be private:
only the two lists of centers and weights need be public.  For boxes, choose
a rational point in each box as center; a proved diameter bound gives a
radius bound.  One may use sharper midpoint radii when available.

These are explicit completeness and multiplicity assumptions, not consequences
of the geometry of the boxes alone.  For polynomial roots they require a
root-counting, factor, or common-refinement theorem.  In the arithmetic
implementation, finite incidence tables, count equations and distance
certificates replace the displayed metric notation.  If actual roots are
represented by names, their analytic interpretation is a separate adapter;
the finite matching program does not inspect those names.

Expand the first family into the $N$ slots
\[
       A^*=\{(i,u):i<m,\ 0\leq u<r_i\},
\]
ordered by $i$ and then $u$, and expand the second family similarly.  A slot
has its cluster's center.  Put $\delta=\epsilon_A+\epsilon_B$ and join
$(i,u)$ to $(j,v)$ when $d(a_i,b_j)\leq\delta$.  Rational $\ell^1$ or
$\ell^\infty$ distances make this a decidable finite graph, without a
square-root operation.  For a general metric interpretation, a supplied
finite adjacency relation suffices if every pair of clusters containing a
common occurrence is an edge and every edge satisfies the displayed
distance bound.

\begin{lemma}[Coverage implies Hall's inequalities]
\label{lem:coverage-hall}
For every set $S$ of slots on the first side, its neighbor set satisfies
$|S|\leq|\mathcal N(S)|$.
\end{lemma}
\begin{proof}
Let $I$ consist of those cluster indices having a slot in $S$, and let $J$
be the cluster indices adjacent to at least one index in $I$.  The triangle
inequality shows that $\alpha(t)\in I$ implies $\beta(t)\in J$:
\[
 d(a_{\alpha(t)},b_{\beta(t)})
 \leq d(a_{\alpha(t)},x_t)+d(x_t,b_{\beta(t)})\leq\delta.
\]
Counting the indicated finite sets gives
\[
 |S|\leq\sum_{i\in I}r_i
   =|\alpha^{-1}(I)|
   \leq|\beta^{-1}(J)|
   =\sum_{j\in J}s_j=|\mathcal N(S)|.
\]
The last equality holds because adjacency depends on cluster indices, not on
which copy of a center is used.  These are bounded finite sums.  In
particular, the argument does not choose a matching between occurrences.
\end{proof}

\begin{theorem}[Bounded constructive matching]
\label{thm:bounded-matching}
There is a primitive recursive algorithm which, on an $N$ by $N$ Boolean
bipartite graph, returns either a perfect matching or a set $S$ with
$|\mathcal N(S)|<|S|$.  Its correctness has a proof by quantifier-free
induction after bounded-formula reflection.  Hence the two enclosures above
have an explicitly computed permutation matching with distance at most
$\epsilon_A+\epsilon_B$ at every matched pair.
\end{theorem}
\begin{proof}
Maintain a matching covering the first $k$ left vertices, initially $k=0$.
To add left vertex $k$, search for an augmenting path.  Mark $k$ on the left.
From a marked left vertex mark all its adjacent right vertices; from a marked
right vertex already in the matching mark its matched left partner.  Record
a parent and a discovery level each time a previously unmarked vertex is
marked.  There are at most $2N$ vertices, so $2N$ rounds suffice to reach
closure; a queue implementation examines at most $N$ left vertices.

If a free right vertex is found, follow the parent records back to $k$.
Discovery levels strictly decrease, so this is a path of fewer than $2N$
edges.  Flip matched and unmatched edges along it.  Every interior vertex
keeps one matched incident edge, and the endpoints acquire one, giving a
matching of size $k+1$ using only graph edges.

If no free right vertex is reached, let $S$ and $T$ be the reached left and
right sets at closure.  Every neighbor of $S$ is in $T$, and each member of
$T$ was reached along an edge from $S$, so $T=\mathcal N(S)$.  Every vertex
of $T$ is matched.  Its partner belongs to $S\setminus\{k\}$.  Conversely,
every reached left vertex other than $k$ was reached through its matched
right partner.  Thus the old matching restricts to a bijection
$T\leftrightarrow S\setminus\{k\}$, and
\[
              |\mathcal N(S)|=|T|=|S|-1.
\]
This is a displayed Hall counterexample.  Therefore Hall's inequalities rule
out failure at each of the $N$ stages.  The final size-$N$ matching is a
permutation because both sides have exactly $N$ vertices.  When $N=0$, the
empty matching is returned immediately.

All scans have the indicated bounds.  Partial matchings, marks, parents,
levels and the stage counter are finite arrays.  Their invariants assert
only bounded index conditions, partial-inverse equations and Boolean edge
tests; each is reflected to one quantifier-free formula.  The outer stage
induction and the bounded search inductions therefore stay within iPRA.
Applying Lemma~\ref{lem:coverage-hall} eliminates the failure branch and gives
the claimed distance bound.
\end{proof}

The algorithm can always inspect vertices in numerical order, so its output
is determined by the \emph{public} adjacency table.  It does not expose a
particular coverage witness $\alpha$ or $\beta$.  This does not assert that
the same permutation is selected when a different adjacency table is used:
what is invariant is the proved matching bound and every justified symmetric
observation.  Private coverage witnesses may be sealed by existential
elimination after the matching theorem is established.

\subsection{A precise observation modulus}

For two length-$N$ lists of points define the matching distance by the finite
minimum
\[
 d_{\mathrm{match}}(a,b)
   =\min_{\sigma\in\mathfrak S_N}\max_{i<N}d(a_i,b_{\sigma(i)}).
\]
The empty-list distance is zero.  This is a finite rational computation when
the underlying distances are rational.  Theorem~\ref{thm:bounded-matching}
proves $d_{\mathrm{match}}(A^*,B^*)\leq\delta$.

Suppose $F$ is invariant under permutations, and that two ordered lists whose
coordinatewise distances are at most $\eta$ have outputs at distance at most
$\omega(\eta)$.  Reorder the second list by the constructed matching and
apply that inequality.  Permutation invariance then gives
\[
                      d(F(A^*),F(B^*))\leq\omega(\delta).
\]
Both invariance and the displayed modulus must be proved for the actual
observation program.  A phrase such as ``stable observation'' does not
replace either obligation.

\begin{proposition}[Polynomial observations]
\label{prop:cluster-polynomial-observations}
Use the submultiplicative norm $\|z\|_1=|\Re z|+|\Im z|$ on Gaussian
rationals, and suppose all centers have norm at most $R$.  If a matching has
error at most $\delta$, then, for $k\geq1$,
\[
 \left\|\sum_{i<N}a_i^k-\sum_{i<N}b_i^k\right\|_1
      \leq NkR^{k-1}\delta.
\]
For $1\leq k\leq N$, the elementary symmetric polynomials satisfy
\[
          \|e_k(a)-e_k(b)\|_1
             \leq \binom{N}{k}kR^{k-1}\delta.
\]
Thus every coefficient of the associated monic product polynomial has an
explicit finite modulus, including at collisions.
\end{proposition}
\begin{proof}
Reorder $b$ by the matching.  The identity
\[
       u^k-v^k=(u-v)\sum_{j=0}^{k-1}u^jv^{k-1-j}
\]
and submultiplicativity give $\|u^k-v^k\|_1\leq kR^{k-1}\delta$.
Summing proves the first assertion.  For each $k$-element subset, subtract
the two corresponding products by changing one factor at a time.  The
telescoping sum has $k$ terms, each bounded by $R^{k-1}\delta$.
There are $\binom{N}{k}$ such subsets, giving the second assertion.
The cases $k=0$ and $k>N$ are constant observations.  No enumeration of
continuous branches is used; every product and sum is finite.
\end{proof}

\subsection{What component balancing alone proves}

Sometimes a verifier supplies only a bipartite graph on the weighted
clusters, an upper bound $\delta$ on each edge length, and equal total red
and blue multiplicity in every connected component.  This weaker
information also gives a matching, with a different quantitative bound.

\begin{proposition}[Component transport]
\label{prop:component-transport}
In a component having $v$ vertices, there is an explicit multiplicity
matching whose every paired distance is at most $(v-1)\delta$.  For $N>0$
and positive weights, the uniform bound $(2N-1)\delta$ is valid.
\end{proposition}
\begin{proof}
List the expanded red and blue slots of the component in numerical order.
Equal total weights give equally long lists; pair equal ranks.  A shortest
edge path between the centers of any paired slots has no repeated vertex,
since deleting the part between two equal vertices shortens it.  It has at
most $v-1$ edges.  The triangle inequality gives $(v-1)\delta$.
Breadth-first search constructs such paths with at most the total number of
vertices as fuel.  Positive weights imply at most $N$ red and at most $N$
blue vertices, giving the uniform bound.  All components and rank selections
are bounded finite computations.
\end{proof}

The factor cannot simply be dropped.  On the real line take red centers
\[
                   0,\quad 1/10,\quad 2
\]
and blue centers
\[
                   1,\quad 29/10,\quad 3,
\]
all of weight one, with edges at distance at most $1$.  The graph is
connected and balanced.  The first two red centers have just the blue center
$1$ as a neighbor, so they violate Hall's inequality.  In any matching one
of these red centers must be paired with $29/10$ or $3$, at distance at
least $14/5$.  Pairing in increasing order attains $14/5$, so the exact
matching distance is $14/5>1$.  This is a counterexample to inferring a
one-edge matching bound from component balance, not to
Theorem~\ref{thm:bounded-matching}, whose complete-coverage assumptions imply
Hall's inequalities and exclude this graph.

For genuine complete covers, component balance itself also follows from the
occurrence witnesses: an occurrence's two assigned cluster vertices are
adjacent, hence in the same component.  The set of occurrences assigned to
red vertices of a component is exactly the set assigned to its blue
vertices.  Counting proves balance.  This supplies the common-coarsening
argument with its precise finite content.

\subsection{Correcting complex labels}

The audit's counterexample is decisive~\cite{repo-fiber}.  For $z^2+1$, take
arbitrarily small rational $h>0$.  Boxes centered at $(h,1)$ and $(-h,-1)$,
each of coordinate half-width $2h$, enclose $i$ and $-i$.  Boxes centered at
$(-h,1)$ and $(h,-1)$ enclose the same roots at the same accuracy.  They
remain separated by a positive metric gap as $h\to0$, but lexicographic
ordering of centers or lower corners reverses.  Thus no reduction in mesh
relative to the metric gap repairs that ordering rule.

There is a simple positive replacement which names exactly the extra data
needed.

\begin{theorem}[Labels relative to stable locators]
\label{thm:stable-locators}
Let $c_0,\ldots,c_{N-1}$ be fixed rational anchors.  Suppose each anchor
has exactly one occurrence within distance $r$, these assignments form a
complete cover, and
\[
                 d(c_i,c_j)\geq2r+\gamma
                 \quad(i\ne j),\qquad \gamma>0.
\]
Let a singleton enclosure of the same occurrences have error
$\epsilon\leq\gamma/4$.  Assign each approximate center to its nearest
anchor, using a bounded minimum scan.  This gives a bijection with the
anchor indices, independent of the singleton-enclosure certificate.
Samples carrying the same label at errors $\epsilon$ and $\epsilon'$ are
within $\epsilon+\epsilon'$ of one another.
\end{theorem}
\begin{proof}
If $x$ is the occurrence assigned to $c_j$ and $a$ approximates $x$, then
\[
                     d(a,c_j)\leq r+\epsilon.
\]
For $i\ne j$, two triangle inequalities give
\[
 d(a,c_i)\geq d(c_i,c_j)-d(c_j,x)-d(x,a)
             \geq r+\gamma-\epsilon.
\]
The second lower bound exceeds the first upper bound by at least
$\gamma-2\epsilon\geq\gamma/2>0$.  Therefore $c_j$ is the unique nearest
anchor; the numerical tie-breaking rule is never used.  The two complete
singleton covers contain exactly one occurrence for each anchor, so every
label is used once.  Two samples with label $j$ approximate that same
occurrence, giving the final inequality by the triangle inequality.
\end{proof}

The labels are relative to the fixed locator certificate.  Approximate
coordinates need not be identical across certificates; the theorem proves
that their labels refer to the same occurrence and that their difference has
the displayed bound.  Taking a dyadic error schedule yields regular names
without any global complex-root ordering.  For real roots, an ordered
interval certificate with positive order gaps gives the familiar simpler
variant.  A certified gap between real parts gives another valid ordering
guard for complex roots.

\subsection{Concrete Lean and iPRA targets}

The target has numerical dimension and graph entries inside the object
language.  Encode an $N$ by $N$ graph by bits of $e<2^{N^2}$; a subset is a
mask $s<2^N$.  Its neighbor mask has bit $j$ equal to the bounded disjunction
of $\operatorname{bit}(s,i)\wedge\operatorname{edge}(e,i,j)$ over $i<N$.
Population count is a bounded sum.  Encode a permutation by $N$ base-$(N+1)$
digits and define
\[
\begin{split}
\operatorname{Match}(N,e,p)\;:\!\!\iff\;&
  (\forall i<N\; p_i<N\ \wedge\ \operatorname{edge}(e,i,p_i)=1)\\
 &{}\wedge(\forall i,j<N\;p_i=p_j\longrightarrow i=j),\\
\operatorname{Defect}(N,e,s)\;:\!\!\iff\;&
 s<2^N\ \wedge\operatorname{count}(\operatorname{neigh}(e,s))
                  <\operatorname{count}(s).
\end{split}
\]
Every displayed quantifier is bounded.  Implement the algorithm by fixed
primitive recursion over its stage, marking and path-flipping bounds.  Its
two output projections $p(N,e)$ and $s(N,e)$ should satisfy the ordinary
closed object theorem
\[
 \forall N,e\;\bigl(e<2^{N^2}\longrightarrow
    (\operatorname{Match}(N,e,p(N,e))
       \vee\operatorname{Defect}(N,e,s(N,e)))\bigr).
\]
The decoder totalizes malformed graphs; the shape guard is retained in the
correctness theorem.  A separate object theorem converts the reflected
coverage/count formulas to the negation of every defect.  Its composition
with the algorithm gives the cluster matching endpoint, with no global
choice principle.

The Lean implementation order is: (1) actual graph/array programs and
their shape equations; (2) matching and alternating-search invariants;
(3) bounded matching-or-defect; (4) finite count monotonicity for the
coverage-to-Hall adapter; (5) actual rational distance and observation
programs; (6) the locator scan and its strict comparison proof.  Reuse
\nolinkurl{Foundations/BoundedSyntax}, \nolinkurl{Foundations/GuardedTrace},
\nolinkurl{Core/SearchParameters}, bounded sums and the rational arithmetic
modules.  The final declarations must return
\texttt{Prf PRAAx []} of these explicit formulas.  A host Lean proof of
Hall's theorem by itself would not establish that endpoint.  The supplied
reference program and finite exhaustive checks are executable evidence for
the algorithm, not a completed kernel proof.
